\chapter{Conclusion} \label{cha:conclusion}
The \textbf{Conclusion} is the final chapter of your thesis. It offers a concise summary of your study, restates your key findings, and underscores their importance. It should leave the reader with a clear understanding of what your work has accomplished—and why it matters.

Unlike the Discussion chapter, which interprets and analyzes the results in depth, the Conclusion provides a high-level synthesis and closure.

\subsection*{Purpose of the Chapter}

The Conclusion should:
\begin{itemize}
  \item Summarize the research aim and main results
  \item Highlight the significance of your findings
  \item Reflect briefly on methodological or theoretical contributions
  \item Suggest directions for future research
\end{itemize}

\subsection*{Suggested Structure}

\begin{enumerate}
  \item \textbf{Restate the Research Purpose} \\
  Begin by reminding the reader of the central research aim and the problem you set out to address.

  \item \textbf{Summary of Main Findings} \\
  Provide a short and focused summary of your most important results. Avoid technical detail—this is a synthesis, not a reanalysis. Structure this part in line with your research questions or study objectives.

  \item \textbf{Contribution of the Thesis} \\
  Emphasize what your thesis contributes to the field. This might be:
  \begin{itemize}
    \item A new empirical finding
    \item A novel application of an existing method
    \item A tool or dataset you created or evaluated
    \item A methodological insight or theoretical clarification
  \end{itemize}

  \item \textbf{Limitations (Optional)} \\
  Briefly acknowledge any important limitations that shape how your conclusions should be interpreted.

  \item \textbf{Outlook / Future Work} \\
  Conclude by suggesting how future research could build on your work. Keep this realistic and grounded—your goal is to show the next logical step, not to promise breakthroughs.
\end{enumerate}

\subsection*{Best Practices}

\begin{itemize}
  \item Keep the conclusion concise—typically 1–2 pages.
  \item Avoid introducing new data, references, or analysis.
  \item Use clear, confident language without overstating your results.
  \item End on a forward-looking note, showing how your work fits into a larger research landscape.
\end{itemize}